% ---------------------------------------------------------------------------
% Пример второй главы / Example chapter 2.
% Демонстрирует систему уравнений с подбрекетами (\myunderbrace), таблицу и
% рисунок. Математические команды определены в Dissertation/userpackages.tex.
% Demonstrates a system of equations with under-braces (\myunderbrace), a table
% and a figure. The math commands are defined in Dissertation/userpackages.tex.
% ---------------------------------------------------------------------------
\chapter{\fixme{Mathematical model}}
\label{ch:model}

\fixme{Introduce the mathematical model of the chapter.}

\section{Governing equations}
\label{sec:governing-equations}

Considering the modelling assumptions, the conservation equations for mass in the
gas and metal phases, together with momentum and energy, are formulated as
follows:
\begin{gather}
  \pdv{\alpha}{t} + \div (\alpha \bv) = \div \bv,
  \label{eq:advection_VOF}\\
  \div \bv = - \frac{(\rho_\liq - \rho_\sol)D_t \phi
    + \phi \dv{\rho_\liq}{T} D_t T}{\phi \rho_\liq + (1-\phi) \rho_\sol},
  \label{eq:mass_conservation}\\
  \begin{multlined}
    \pdv{\rho\bv}{t}
    + \div (\rho \bv \otimes \bv)
    = - \grad p
    + \grad \vdot \btau
    + \myunderbrace{\rho \bg}{gravity}
    - \myunderbrace{K\alpha\frac{(\alpha-\phi)^2}{\phi^2 + \epsilon}\bv}{Brinkman penalization} \\
    - \myunderbrace{\gamma\qty(\div\bn_\alpha)\grad\alpha}{surface tension}
    + \myunderbrace{\dv{\gamma}{T}\Bigl(\bn_\alpha\cp\grad{T}\cp\bn_\alpha\Bigr)|\grad \alpha|}{Marangoni force},
    \label{eq:momentum_conservation}
  \end{multlined}\\
  \pdv{\rho h}{t}
  + \div (\rho \bv h)
  = \div (k \grad{T})
  + \myunderbrace{a_\laser I_{\laser} \max(0, \bn_\laser \vdot \grad \alpha)}{laser heat source},
  \label{eq:h_eq}
\end{gather}
where $\bv$ is the velocity of the mixture, $p$ is the dynamic pressure,
$\btau = \mu [\grad \bv + (\grad \bv)^\tran] - \frac{2}{3}\mu(\div \bv)\bI$ is the
viscous stress tensor, $\bg$ is the gravitational acceleration, $\gamma$ is the
surface tension, $k$ is the thermal conductivity, $I_{\laser}$ is the laser beam
intensity with absorption coefficient $a_\laser$, $\bn_\laser$ is the laser beam
direction, and $D_t \equiv \pdv*{t} + \bv\vdot\grad$ is the material derivative.
The momentum balance~\eqref{eq:momentum_conservation} is written with each force
term highlighted by an under-brace.

The thermophysical properties used in the model are listed in
Table~\ref{tab:properties}.

\begin{table}[htbp]
  \caption{\fixme{Thermophysical properties of the metal and the auxiliary gas
  phase.}}
  \centering
  \renewcommand{\arraystretch}{1.2}
  \small
  \begin{tabular}{cccc}
    \hline
    Phase & Physical property & Value & Unit \\ \hline
    \multirow{5}{*}{Liquid metal}
        & $\rho_{\liq\melt}$ & \num{2475}   & \mysi{\kg\per\cubic\m}   \\
        & $\beta$            & \num{-1.4e-4}& \mysi{\per\K}            \\
        & $k_\liq$           & \num{91}     & \mysi{\W\per\m\per\K}    \\
        & $c_{p\liq}$        & \num{1042}   & \mysi{\J\per\kg\per\K}   \\
        & $\mu_\liq$         & \num{4.0e-5} & \mysi{\Pa\s}            \\ \hline
    \multirow{3}{*}{Interface}
        & $\gamma$           & \num{8.7e-1} & \mysi{\N\per\m}          \\
        & $T_\melt$          & \num{936.6}  & \mysi{\K}               \\
        & $h_\fus$           & \num{3.8e5}  & \mysi{\J\per\kg}         \\ \hline
    \multirow{3}{*}{Solid metal}
        & $\rho_\sol$        & \num{2700}   & \mysi{\kg\per\cubic\m}   \\
        & $k_\sol$           & \num{211}    & \mysi{\W\per\m\per\K}    \\
        & $c_{p\sol}$        & \num{910}    & \mysi{\J\per\kg\per\K}   \\ \hline
    \multirow{4}{*}{Auxiliary gas phase}
        & $\rho_\gas$        & \num{4.0e-1} & \mysi{\kg\per\cubic\m}   \\
        & $\mu_\gas$         & \num{4.0e-5} & \mysi{\Pa\s}            \\
        & $k_\gas$           & \num{6.1e-3} & \mysi{\W\per\m\per\K}    \\
        & $c_{p\gas}$        & \num{e3}     & \mysi{\J\per\kg\per\K}   \\ \hline
  \end{tabular}
  \label{tab:properties}
\end{table}

\section{Numerical illustration}
\label{sec:numerical-illustration}

A diagram drawn directly in \LaTeX{} with \textsf{TikZ} (its axes labelled with
math) is shown in Figure~\ref{fig:enthalpy-explanation}: the specific enthalpy
$h$ along the coordinate $x$ between the nodes $x_l$ and $x_r$.

\begin{figure}[htbp]
    \centering
    \begin{tikzpicture}[>=latex', scale=2.0]
        \def\xmax{3.4}
        \def\xstart{0.5}
        \def\ystart{0.5}

        \coordinate (A) at (\xstart,\ystart+0.25);
        \coordinate (L) at (\xstart+0.5,\ystart+1.25/2);
        \coordinate (LX) at (\xstart+0.5,0);
        \coordinate (LH) at (0,\ystart+1.25/2);
        \coordinate (B) at (\xstart+1,\ystart+1);
        \coordinate (BX) at (\xstart+1,0);
        \coordinate (BH) at (0,\ystart+1);
        \coordinate (M) at (\xstart+1.5,\ystart+1.5);
        \coordinate (MX) at (\xstart+1.5,0);
        \coordinate (MH) at (0,\ystart+1.5);
        \coordinate (C) at (\xstart+2,\ystart+2);
        \coordinate (CX) at (\xstart+2,0);
        \coordinate (CH) at (0,\ystart+2);
        \coordinate (R) at (\xstart+5/2*1.1,\ystart+4.5/2*1.05);
        \coordinate (RX) at (\xstart+5/2*1.1,0);
        \coordinate (RH) at (0,\ystart+4.5/2*1.05);
        \coordinate (D) at (\xstart+3,\ystart+2.5);

        % AXES --- labelled in LaTeX math
        \draw[->,thick] (0,0) -- (0,\xmax+0.2) node[left=0] {$h$};
        \draw[->,thick] (0,0) -- (\xmax+0.2,0) node[below=0] {$x$};

        % Lines
        \draw[thick,dashed] (A) -- (B);
        \draw[thick]        (B) -- (C);
        \draw[thick,dashed] (C) -- (D);
        \draw[thick,dashed] (L) -- (LX);
        \draw[thick,dashed] (R) -- (RX);
        \draw[thick,dashed] (B) -- (BX);
        \draw[thick,dashed] (B) -- (BH);
        \draw[thick,dashed] (M) -- (MX);
        \draw[thick,dashed] (M) -- (MH);
        \draw[thick,dashed] (C) -- (CX);
        \draw[thick,dashed] (C) -- (CH);
        \draw[thick,dashed] (L) -- (LH);
        \draw[thick,dashed] (R) -- (RH);

        \fill[black] (B) circle(0.04);
        \fill[black] (BH) circle(0.04) node[fill=white,inner sep=0.5,left=0.5] {$h_1$};
        \fill[black] (L) circle(0.04);
        \node[inner sep=0.5, shift={(-0.35, 0.35)}] at (L) {$\phi = 0$};
        \fill[black] (LX) circle(0.04) node[fill=white,inner sep=0.5,below=0.5] {$x_l$};
        \fill[black] (LH) circle(0.04) node[fill=white,inner sep=0.5,left=0.5] {$h_l$};
        \fill[black] (R) circle(0.04);
        \node[inner sep=0.5, shift={(-0.35, 0.35)}] at (R) {$\phi = 1$};
        \fill[black] (M) circle(0.04);
        \fill[black] (MH) circle(0.04) node[fill=white,inner sep=0.5,left=0.5] {$h_m$};
        \fill[black] (RX) circle(0.04) node[fill=white,inner sep=0.5,below=0.5] {$x_r$};
        \fill[black] (RH) circle(0.04) node[fill=white,inner sep=0.5,left=0.5] {$h_r$};
        \fill[black] (C) circle(0.04);
        \fill[black] (CH) circle(0.04) node[fill=white,inner sep=0.5,left=0.5] {$h_2$};
    \end{tikzpicture}
    \caption{\fixme{Specific enthalpy $h$ at the nodes $x_l$, $x_r$ and in the
    region between them, illustrating the liquid volume fraction $\phi$.}}
    \label{fig:enthalpy-explanation}
\end{figure}

\section{Solidification shrinkage}
\label{sec:solidification}

\fixme{Discuss the effect of solidification shrinkage on the free surface.}
Figure~\ref{fig:solidification-shrinkage} shows the melt-pool shape at different
time instants during solidification for several density models.

\begin{figure}[htbp]
  \centering
  \includegraphics[width=0.9\linewidth]{pictures/ver_sol_diff_density}
  \caption{\fixme{Solidification of the melt pool at different time instants for
  several density models, illustrating volumetric shrinkage.}}
  \label{fig:solidification-shrinkage}
\end{figure}
